\scanpage{78}
\unsure{On arrive donc encore} $H_A^1(G)=G(K)/G(A)$, qui est nul si $\dim A>1$. \unsure{+|et}

[Supposons \uline{$\dim A=1$}. On aura donc $H_A^0(G)=0$
\diagram{\begin{tikzcd}[ampersand replacement=\&,column sep=small]
{}\arrow[r,no head]\&H^1(A,G)\arrow[r,no head]\&H^1(K,G)\arrow[r,"\unsure{\partial|\delta}"]\&H^2(A\unsure{\longrightarrow|\leadsto}K,G)\arrow[r]\&H^2(A,G)\arrow[r,no head]\&H^2(K,G)\arrow[d,"\partial"]\\
\&\&\&\&\&H^3(A\longrightarrow K,G)
\end{tikzcd}}
\unsure{donc en dim} $\unsure{1}$ les $H^i(A\longrightarrow K,G)=H_A^i(G)$ \unsure{sont extensions} d'un noyau par un conoyau de
\diagram{\begin{tikzcd}[ampersand replacement=\&]
H^i(A,G)\arrow[r]\arrow[d,equal]\&H^i(K,G)\\
H^i(L,G)\&
\end{tikzcd}}
\deleted{Si $G$ est}
\unsure{Noter} que si $\unsure{L|\mathscr L}$ est le plus \unsure{gd} ss-gpe linéaire connexe de $G$, alors
\[
\left.\begin{aligned}
H^i(A,G)&=H^i(A,G/L)\\
H^i(K,G)&=H^i(K,G/L)
\end{aligned}\;\right|\quad i\geq1
\]
donc on peut supposer $L=0$]

\diagram{\begin{tikzcd}[ampersand replacement=\&,column sep=large,row sep=large]
\unsure{e|\epsilon|0}\arrow[d,no head]\&\&\\
\pi_1^{\deleted{\unsure{a}}}(L)=\pi_1(A)\&\pi'\&{\unsure{B|\mathbb B}}_{\pi'}\\
\pi_1(K)\arrow[u]\&\pi\&B_\pi\arrow[u]\\
\pi_1(A\longrightarrow K)\arrow[u]\&\pi''\&B_{\pi''}\arrow[u]\\
e\arrow[u]\&\&\unsure{\pi'|\pi_1|\Pi}\arrow[u]
\end{tikzcd}}
% archive sheet number: 77
